\runninghead{The Seed That Is Crushed}
\properlane{crushed-seed}{The Least of All Seeds: Strength Given Out When the Seed Is Crushed}

The parable says two things of the mustard seed that do not sit easily
together. It is \latin{mínimum \dots\ ómnibus semínibus}, ``the least \dots\ of
all seeds,'' and when it has grown it is greater than all the garden herbs
\latin{et fit arbor}, ``and becometh a tree,'' with birds lodging in its
branches. Mustard is an herb, and the parable itself counts it among the
herbs before it calls it a tree. The Fathers ask what the small thing is and
what makes it great. St Hilary of Poitiers and St Ambrose answer from a
property of the seed: mustard is sharp, and its sharpness comes out when it is
crushed. The seed for them is the Lord, given up to death, and, for Ambrose,
also the faith that believes in him. This interpretation sets beside them the
church of the Epistle, which received the word in much tribulation before its
faith was spoken of everywhere.

\subsection{The Lord compared himself to the mustard seed}
\label{sec:seed-hilary}

Hilary opens his comment with the identification: \latin{Grano sinapis se
ipsum Dominus comparavit, acri maxime, et omnium seminum minimo, cujus virtus
ac potestas tribulationibus et pressuris accenditur}, ``The Lord compared
himself to a grain of mustard, the sharpest and the least of all seeds, whose
strength and power is kindled by tribulations and pressures''
(\work{In Matthaeum}~13.4, PL~9 col.~994). He reads the parable through the
Passion and the mission that follows it. The seed is sown in the field when the Lord,
\latin{a populo comprehensus et traditus morti}, ``seized by the people and
handed over to death,'' was buried, so that his burial was a kind of sowing.
The herbs that the grown plant surpasses are the prophets: their preaching
had been given to a sick Israel \latin{oleris \dots\ vice}, as vegetables are
given to the sick, and the Lord grows beyond \latin{universam prophetarum
gloriam}, ``the whole glory of the prophets.'' The branches are the apostles,
and the birds are the nations, tossed about by the devil's gusts, who
\latin{in spem vitae advolabunt}, ``will fly to them in hope of life.''

The seed's smallness Hilary simply states, without giving a cause, and he
does not take up how an herb can become a tree; Jerome and Bede, below,
compare the herb and the tree.

\subsection{Ground, pressed and sown: Ambrose on faith and on Christ}
\label{sec:seed-ambrose}

Ambrose expounds the same parable as Luke gives it, where the man sows the
seed in his garden (Lk 13:18--19), and quotes Matthew's wording of it. He
begins from the other saying in which mustard appears, which he quotes as
\latin{Si habueritis fidem sicut granum sinapis, dicetis huic monti: Tollere,
et iactare in mare}, ``If you have faith like a grain of mustard seed, you
will say to this mountain, Be lifted up and cast into the sea'' (Mt 17:19, where
the mountain is told \latin{Transi hinc illuc}, ``Remove from hence hither'';
the command Ambrose quotes is that of the mountain saying at Mt 21:21 and
Mk 11:23). Since both the kingdom and faith are compared to the mustard seed,
he concludes \latin{fides utique regnum coelorum est}, ``faith, then, is the
kingdom of heaven'' (\work{Expositio in Lucam}~VII.176--177). He then asks
what the seed's nature teaches: \latin{Granum eius certe res est vilis et
simplex: si teri coeperit, vim suam fundit. Etiam fides primo simplex videtur:
sed si teratur adversis, gratiam suae virtutis effundit}, ``Its grain is
indeed a cheap and simple thing: if it begins to be ground, it pours out its
strength. Faith too at first seems simple: but if it is ground by adversity,
it pours out the grace of its power,'' and fills with its scent those who hear
or read of it (VII.178).

His first example is from his own church. The martyrs Felix, Nabor and
Victor, whom he calls \latin{martyres nostri}, ``our martyrs,''
\latin{habebant odorem fidei, sed latebat}, ``had the scent of faith, but it
was hidden.'' Persecution came; they laid down their arms and bent their
necks, and \latin{contriti gladio per totius terminos mundi gratiam sui
sparsere martyrii}, ``crushed by the sword, they spread the grace of their
martyrdom to the ends of the whole world,'' so that the psalm can rightly be
said of them, \latin{In omnem terram exivit sonus eorum}, ``Their sound went
out into all the earth'' (Ps 18:5; VII.178).

Then he turns to the Lord. \latin{Ipse Dominus granum est sinapis}, ``The Lord
himself is the grain of mustard.'' He chose to be ground and pressed and sown:
\latin{in horto enim Christus et captus et sepultus: in horto crevit, ubi etiam
resurrexit, et factus est arbor}, ``for Christ was both taken and buried in a
garden; in a garden he grew, where he also rose, and became a tree''
(VII.179). The hearer is told to do the same with his own faith: \latin{semina
in horto tuo Christum}, ``sow Christ in your own garden'' (VII.180), and
\latin{Premitur fides, cum crucifixum credimus Christum}, ``faith is pressed
when we believe in Christ crucified,'' as Paul pressed it when he came to the
Corinthians \latin{non cum eminentia sermonis}, ``not with eminence of
speech,'' and judged that he knew nothing among them \latin{nisi Iesum
Christum, et hunc crucifixum}, ``but Jesus Christ, and him crucified'' (1~Cor
2:1--2; VII.181). To
the question how Christ can be the least of seeds, Ambrose answers that he is
least \latin{per speciem}, in appearance, not in nature: he \latin{non
habebat speciem neque decorem}, ``had no form nor comeliness'' (Is 53:2), and
he is \latin{speciosus forma prae filiis hominum}, ``beautiful above the sons
of men'' (Ps 44:3);
and he surpassed \latin{omnem gloriam prophetarum, quos infirmus Israel velut
olera manducaverat}, ``all the glory of the prophets, whom a weak Israel had
eaten like herbs'' (VII.183), which is Hilary's thought taken up without his
name.

Ambrose holds the two identifications together: the seed is faith, because
faith gives its strength under adversity, and it is Christ, because he was
ground and buried and rose. St Peter Chrysologus, preaching on Luke's parable, says that
Christ, sent into the garden of the Virgin's body, \latin{orbe toto crucis in
arborem crevit}, ``grew on the cross into a tree over the whole world,'' and
gave his savour \latin{cum teritur passione}, ``when he was ground in the
Passion.'' He gives the reason from the seed: \latin{sicut in integritate
grani sinapis virtus habetur occulta, apparet vero virtus ejus
vehementissime si teratur; ita Christus corpore teri voluit, qui suam noluit
latere virtutem}, ``as the power of the mustard grain is hidden while it is
whole, but shows most strongly if it is ground, so Christ willed to be ground
in his body, who did not will that his power should stay hidden''
(\work{Sermo}~98). St Bede, after giving Jerome's and Augustine's readings of
Luke's parable, adds as a further sense \latin{ipsa dominicae incarnationis
humilitas}, ``the very humility of the Lord's incarnation'': Joseph took the
crucified Saviour's body and buried it in a garden; it rose and ascended, and
became a tree whose branches are the preachers sent through the world
(\work{In Lucam}~IV, on 13:19). Ambrose had already taken the burial in the
garden for the sowing (VII.179).

\subsection{Other answers: the preaching, the apostles, the heat of faith}
\label{sec:seed-others}

The Fathers do not agree on what the seed is, and their other answers place
the smallness elsewhere. St Jerome takes the kingdom to be the preaching of
the gospel, and the preaching is \latin{minima \dots\ omnibus disciplinis},
``the least of all teachings,'' because at first it seems untrue: it preaches
\latin{hominem Deum, Deum mortuum, et scandalum crucis}, ``a man who is God,
God dead, and the scandal of the cross.'' Set beside the philosophers' books
and their eloquence it looks small; but their doctrines grow quickly into
herbs that wither, while the preaching grows into a tree (\work{In
Matthaeum}~II, on 13:31--32, PL~26 col.~90). St Thomas Aquinas follows him:
the doctrine seemed least because it preached God \latin{passum,
crucifixum}, ``suffering, crucified,'' and it grew greater than other
doctrines \latin{in soliditate, in generalitate, et in utilitate}, ``in
solidity, in reach and in usefulness'' (\work{Super Evangelium Matthaei}
c.~13). St Bede, on Mark's parable, gives the tree three advantages over the
herbs, height, breadth and age: it is broad because \latin{per praedicatores
ex sese nascentes, totius mundi terminos occupavit}, ``through the preachers
born from it, it has taken possession of the ends of the whole world,'' and
it lasts because \latin{veritas quam praedicat, nullo unquam fine delebitur},
``the truth it preaches will never be destroyed'' (\work{In Marcum}~I, on
4:30--32). St Augustine reads the seed by its heat: it becomes greater than
all herbs, that is, than the opinions of the sects, \latin{ob fervorem
fidei}, ``because of the fervour of faith'' (\work{Quaestiones
Evangeliorum}~I.11).

Jerome's answer is the nearest to Hilary's and Ambrose's and is still a
different answer. For Jerome the seed is small because of what it says: the
cross is a scandal to those who hear it preached. He does not say that the
preaching gives its strength when it is persecuted. Ambrose's ``faith is
pressed when we believe in Christ crucified'' touches Jerome's point, since
for both the cross is what makes the seed look small; but Ambrose then speaks
of martyrs ground by the sword, and Jerome does not.

St John Chrysostom reads the parable as the Lord's answer to his disciples'
fear. After the parable of the sower, in which three parts of the seed are
lost, he imagines them asking, ``And who, and how many will be the
faithful?''; the Lord answers that fear with the mustard seed, ``signifying
that in any case the gospel shall be spread abroad.'' The smallness is theirs: ``His disciples were weakest of all,
and least of all; but nevertheless, because of the great power that was in
them, It hath been unfolded in every part of the world'' (\work{Hom.\ in
Mt.}~46.2). Aquinas reports this as Chrysostom's reading, the seed being the
apostles, \latin{modici et abiecti}, ``small and despised,'' who did more than
Alexander and the Romans. For Chrysostom, as for Hilary and Ambrose, the power
that makes the small thing great is God's; what is small is, for him, the
disciples, and for Ambrose the Lord's outward appearance.

\subsection{A church that received the word in much tribulation}
\label{sec:seed-thessalonica}

Paul's converts at Thessalonica received the word \latin{in
tribulatióne multa}, ``in much tribulation'' (1~Thess 1:6), and two verses
later he says that from them \latin{diffamátus est sermo Dómini}, ``was spread
abroad the word of the Lord,'' and that their faith had gone out \latin{in omni
loco}, ``in every place'' (1:8). Paul tells of this young church in the order
of Ambrose's martyrs, who were crushed by the sword and then spread the grace
of their martyrdom to the ends of the world (VII.178): the tribulation first,
and then the word sounding out from them everywhere.

Chrysostom's homilies on the letter say why the pressure came and how it was
borne. On
``the work of your faith'' he writes, ``If thou believest, suffer all things;
if thou dost not suffer, thou dost not believe.'' He explains why the
persecution came at once: ``God permitted that persecutions should arise
immediately, that no one might say, that the Preaching was established lightly
or by flattery,'' and so that it might be clear that ``it was not human
persuasion, but the power of God, that persuaded the souls of the believers''
(\work{Hom.\ in 1~Thess.}~1). On ``in much tribulation, with joy of the Holy
Ghost'' he answers the question the phrase raises, how affliction and joy can
meet:
``The affliction is in things bodily, and the joy in things spiritual,'' as the
Three Children in the furnace were refreshed with dew. They became imitators
``of the Lord'' because he too suffered and rejoiced. On verse~8 he notes the
verb: Paul ``has not said, your faith is noised abroad, but `has sounded
out','' as every place near is filled with the sound of a loud trumpet
(\work{Hom.}~2). Aquinas, on the Latin, says that the joy is the Holy
Spirit's and no other's, the Spirit ``who is the love of God, who gives joy to
those who suffer for Christ, because they love him'' (\latin{qui est amor Dei,
qui facit gaudium patientibus propter Christum, quia amant eum};
\work{Super I ad Thessalonicenses} c.~1 l.~1).

Paul also says how the gospel came: \latin{non \dots\ in sermóne tantum, sed et
in virtúte}, ``not \dots\ in word only, but in power also'' (1:5). Aquinas
sets this against sublimity of speech, citing Paul's account of his preaching
at Corinth, \latin{non in persuasibilibus humanae sapientiae verbis, sed in
ostensione spiritus et virtutis}, ``not in the persuasive words of human
wisdom, but in shewing of the Spirit and power'' (1~Cor 2:4). He then gives
two explanations of the power. Either the preaching was confirmed \latin{non
argumentis, sed virtute miraculorum}, ``not by arguments but by the power of
miracles'' (Mk 16:20; Heb 2:4), or Paul showed them \latin{virtuosam vitam},
``a virtuous life'' (\work{Super I ad Thessalonicenses} c.~1 l.~1).

At the opening of the same chapter Paul says that he came to Corinth not
with eminence of speech and judged that he knew nothing among them but Christ
crucified (1~Cor 2:1--2); these are the verses Ambrose cites for faith pressed
by belief in the crucified (VII.181). At Thessalonica, too, he had reasoned
in the synagogue that ``the Christ was to suffer and to rise again from the
dead'' (Acts 17:2--3). On Ambrose's account, then, the Thessalonians' faith
was pressed in what they believed as well as by the persecution that
followed.

Chrysostom draws from the parable a conclusion about the hearers. ``Now if
twelve men leavened the whole world, imagine how great our baseness, in that
when we being so many are not able to amend them that remain.'' He refuses
the excuse that the apostles had miracles: ``It was not the signs that made
them admirable,'' but ``their contempt of wealth, their despising glory, their
freedom from worldly things \dots\ it is their life, so bright on all sides,
which also draws down the grace of the Spirit'' (\work{Hom.\ in Mt.}~46.3). The
Collect asks that those who pray may carry out what pleases God \latin{et
dictis \dots\ et factis}, ``both in word and in deed.'' The deeds it asks for
can be understood as a life of the kind Chrysostom praises. Paul's own pair in
1:5 is word and power. On Aquinas's second explanation that power is the
virtuous life Paul showed the Thessalonians, which is close to what Chrysostom
praises in the twelve; on his first it is the miracles that confirmed the
preaching.

\subsection{The afflicted who sing the chants}
\label{sec:seed-chants}

Three of the chants are spoken to or by people under pressure. The Introit
sings God's words to exiles in Babylon, a small and captive people far from their land: his
thoughts toward them are of peace, and he will bring them back \latin{de
cunctis locis}, ``from all places.'' Jerome's gloss on the clause the
antiphon omits, ``to give you an end and patience,'' speaks of \latin{patientiam
laborum praesentium}, ``patience in the present labours'' (\work{In
Hieremiam}~V, PL~24 col.~893).

The Gradual's psalm (see its setting above) thanks God for past
deliverance and then complains that God has cast his people off and that they
are killed all the day long for his sake (Ps 43:10, 22). Augustine hears in it
``the voice of the Martyrs'' (\work{En.\ in Ps.}~43.1). Chrysostom, on the
Gradual's verses, takes ``all the day'' as the whole of life, in which we glory
in God not only when he helps us but also when he forsakes us, as Paul glories
only in the cross (Gal 6:14; \work{Expositio in Ps.}~43, on vv.~8--9).

The Alleluia and the Offertory sing the cry from the depths. Hilary places the
one who cries in three depths: \latin{pone illum \dots\ in profundo divitiarum
Dei}, ``place him in the depth of the riches of God,'' unsearchable and past
finding out (Rom 11:33); \latin{statue illum de profundo humanarum calamitatum
aestu precantem}, ``set him before you praying from the deep surge of human
calamities''; or \latin{existima illum non labiis, sed secreto, ut jussum est,
cordis orare}, ``think of him as praying not with his lips but in the secret
of the heart, as we are commanded'' (Mt 6:6). The cry, he says, fits
\latin{aut ignorationi, aut passioni, aut orationi}, ``ignorance, or
suffering, or prayer'' (\work{Tractatus in Ps.}~129.2). Those who are pressed
cry from his second depth. Chrysostom says that ``out of the depths''
means from the deepest heart, as the souls of the afflicted pray; prayer of
this kind, rooted deep like a strong tree, stands
against every gust of the devil, while prayer from the lips is overturned
(\work{Expositio in Ps.}~129.1).

The Secret is said over the offering after that cry. The Church holds that
the offering of the Mass is the sacrifice of the cross made present, the death
in which Hilary sees the seed sown. The Secret asks that this offering
\latin{mundet \dots\ et rénovet, gubérnet et prótegat}, cleanse and renew,
govern and protect, those who offer it.

\subsection{Faith that moves the mountain}
\label{sec:seed-communion}

Ambrose began his exposition from the saying that faith as a grain of mustard
seed can move a mountain, and Jerome reads the same saying in his comment on
the parable: \latin{quod Dominus dixit fide fieri, quae grano sinapis
comparetur, hoc Apostolus docet tota fide posse fieri}, ``what the Lord said is
done by faith that is compared to a grain of mustard, the Apostle teaches can
be done by all faith'' (1~Cor 13:2; \work{In Matthaeum}~II, PL~26 col.~91).
Mustard-seed faith, for both, is great faith, the faith that moves mountains.
In Mark the Communion's verse follows directly on a saying about that faith:
``whosoever shall say to this mountain, Be thou removed \dots\ and shall not
stagger in his heart, but believe \dots\ it shall be done unto him'' (Mk
11:23). The Communion's \latin{crédite quia accipiétis}, ``believe that you
shall receive,'' is therefore spoken to the faith that Ambrose and Jerome call
the mustard seed: Mark puts the promise immediately after the mountain, and
both Fathers explain the mustard seed by the saying about the mountain.

St Bede, on Mark's verses, distinguishes those who pray. \latin{Qui perfectam
habet fidem, quae per dilectionem operatur, ille orando, vel etiam jubendo,
potest montes transferre spiritales}, ``he who has perfect faith, which works
through love, can by praying, or even by commanding, move spiritual
mountains.'' His examples are Paul's: the magician Elymas, and the spirit
Paul cast out of the girl at Philippi, \latin{montem videlicet
superbissimum}, ``a most proud mountain''; and he adds that the persecution of
the gentiles that followed at once showed how much fury that mountain, cast
into the sea, brought on Paul (\work{In Marcum}~III, on 11:22--25). The
persecution at Philippi is the one Paul recalls to the Thessalonians: he came
to them ``having suffered many things before and been shamefully treated
\dots\ at Philippi'' (1~Thess 2:2).

Chrysostom also places the promise among dangers. On Matthew's form of it
(Mt 21:21--22), he says that it was given ``that He might train them to feel confidence'': ``Ye also shall do
greater things, if ye are willing to believe and to be confident in prayer,''
said ``so that they might not be afraid and tremble at plots against them,''
and repeated ``to make them cleave to prayer and faith'' (\work{Hom.\ in
Mt.}~67.2).

The Postcommunion asks that those who have been fed may \latin{semper}, always,
long for the things by which they truly live. Bede says of the preaching that
it grows also \latin{in mente cujusque credentis, quia nemo repente fit
perfectus}, ``in the mind of each believer, because no one is made perfect all
at once,'' citing the psalm of the ascent by stages (Ps 83:6--8; \work{In
Lucam}~IV, on 13:19). By Bede's account those who have been fed are still
growing, and the Postcommunion asks that their desire go on.

\begin{fourSenses}
\item[Literal.] A mustard seed, which the parable calls the least of all
seeds, grows past the garden herbs and becomes a tree in which the birds
lodge; Jesus tells this to the crowds by the sea (Mt 13:1--2, 31--32). Paul
writes to a church founded in a few weeks and persecuted at once (Acts
17:1--10) that it received the word in much tribulation and that its faith is spoken of in every place (1~Thess
1:6--8).
\item[Allegorical.] The seed is Christ, seized, put to death and buried, who
rose and filled the world through his apostles (Hilary, \work{In
Matthaeum}~13.4; Ambrose, VII.179--181; Chrysologus, \work{Sermo}~98); for
Bede, the humility of the incarnation, buried in the garden and risen. The
herbs he outgrows are the prophets; the birds are the nations who fly to the
apostles in hope of life (Hilary).
\item[Moral.] Faith gives out its strength when adversity presses it
(Ambrose, VII.178); the believer sows Christ in his own garden (VII.180) and,
believing, suffers all things (Chrysostom, \work{Hom.\ in 1~Thess.}~1). A few
who live the gospel can leaven many, and it was the apostles' life and not
their signs that made them great (Chrysostom, \work{Hom.\ in Mt.}~46.3).
\item[Anagogical.] In the grown tree's branches the one caught up into
paradise, and to be caught up into the clouds, rests \latin{tuta iam sede},
``in a seat now safe'' (2~Cor 12:4; 1~Thess 4:16); the nations gathered from
far off fly into them and flourish \latin{certae stationis iam securitate},
``in the security of a sure station'' (Ambrose, VII.185--186). The nations
fly to the branches in hope of life (Hilary); the tree lasts, because the
truth it preaches will never end (Bede, on Mk 4:30--32).
\end{fourSenses}
